%% Interstitial: Paper 0 -> Paper 1
%% Narrative: the probability exists; now we ask how it relates to
%% the observable experiments that generated it.

Chapter~\ref{chap:paper-zero} established that compatible observable
laws determine a unique $\sigma$-additive probability measure $P$ on
$(\Omega, \sigma(\mathcal{Q}))$.  The measure exists.  It is
canonical.  Its marginals recover the original valuations $\{\ell_Q\}$.

But the construction was existential.  We know $P$ exists and is
unique; we do not yet have a structural account of what $P$ \emph{is}
--- of how it is built from the query system, what properties it
inherits from the $\{\ell_Q\}$, and when two different query systems
determine the same $P$.

\medskip

These are the questions of Chapter~\ref{chap:paper-one}.  Its central
result is an \emph{observational determination theorem}: two
$\sigma$-additive probability measures on $(\Omega, \sigma(\mathcal{Q}))$
that agree on all query marginals are identical.  Probability on the
realization space is entirely determined by observable experiments ---
not partially, not generically, but exactly and completely.

This is the program's first structural theorem, and it has a
philosophical weight beyond its technical content.  It says that there
is no hidden probability: no aspect of $P$ that the observable
interface fails to capture.  The latent space $\Omega$ carries exactly
the structure forced by the queries, and no more.

\medskip

Chapter~\ref{chap:paper-one} also establishes two independent routes
to the global measure, corresponding to two different starting
hypotheses about the marginals: the \emph{realizability route}
(marginals already $\sigma$-additive, realizability condition on
$\Omega$) and the \emph{Prokhorov route} (marginals only finitely
additive, SPUT supplies $\sigma$-additivity).  The relationship between
these routes --- and the conditions under which they coincide ---
clarifies the logical structure of what SP2 required.

\medskip

The passage from Chapter~\ref{chap:paper-zero} to
Chapter~\ref{chap:paper-one} is the passage from \emph{existence} to
\emph{structure}: from knowing that a probability exists to
understanding how it is determined, when it is unique, and what
observing it fully means.
